\scanpage{9}

b) $\tau_i$ est $\delta_i$-exact

c) Pour tout $E\in\mathrm{ob}\,\Delta_f$, le foncteur $\varphi\longmapsto\varphi\circ\tau_i$ de $\underline{\mathrm{Hom}}_\Delta(C_{i+1},E)\longrightarrow\underline{\mathrm{Hom}}_\Delta(C_i,E)$ induit une \uline{équivalence}
\[
\underline{\mathrm{Hom}}_{\Delta,\delta_{i+1}}(C_{i+1},E)\xrightarrow{\approx}\underline{\mathrm{Hom}}_{\Delta,\delta_i}(C_i,E).
\]

\deleted{On pose}

De c) on déduit par récurrence transfinie, pour $E\in\mathrm{ob}\,\Delta_f$
\[
(c^*)\qquad \underline{\mathrm{Hom}}_{\Delta,\delta_i}(C_i,E)\xrightarrow{\approx}\mathrm{Hom}_{\Delta,\delta}(C,E)
\]

et de la définition des $\delta_i$ on déduit par réc. transfinie
\[
\delta_i=\alpha_i(\delta)\qquad(\alpha_i:C\unsure{-|\to}C_i)
\]

Donc \uline{tout revient à prouver que l'on peut trouver un} (petit) \uline{ordinal $\omega$ tel que} $C_\omega\in\mathrm{Ob}\,\Delta_f$, \inserted{\uline{et} $C-C_\omega$ est $\delta$-exact.}

Pour ceci, soit $\pi$ un cardinal majorant strict$^{\mathrm{mt}}$ $\mathrm{Fl}\,R$ et $\mathrm{Fl}\,R'$, \deleted{$\mathrm{Fl}\,\Theta\ \unsure{\mathrm{et}}$} pour les $R,R'$ qui interviennent dans $f$, ainsi que les \inserted{$\mathrm{Fl}\,\Theta$ si $\Theta$ intervient dans $\delta'$ ou $\delta''$.} Soit $\omega$ le plus petit ordinal de cardinal $\pi$. Donc l'ens $S_\omega$ des $j<\omega$ est grand devant $\mathrm{card}\,\mathrm{Fl}\,\unsure{,|{}}R$, $\mathrm{card}\,\mathrm{Fl}\,R'$, $\mathrm{Fl}\,\Theta$.

Mais on voit que $C-C_\omega$ est $\delta$-exact en utilisant b) et le fait que $S_\omega$ grand devant les \inserted{$\mathrm{card}$} $\mathrm{Fl}\,\Theta$. Et on voit que pour $\varphi:R-R'$, $\underline{\mathrm{Hom}}(R',C_\omega)-\underline{\mathrm{Hom}}(R,C_\omega)$ est resp. fidèle ($\varphi$ un $f^{(i)}_\lambda$, avec $i=0,1,2$) ou pl. fidèle $(i=1,2)$, ou équivalence $(\unsure{i|\tau}=2)$, en utilisant respectivement $a_0,a_1,a_2$ plus \unsure{les|le} fait que $S_\omega$ est grand devant \deleted{\unsure{$\mathrm{Fl}$}} $\mathrm{card}\,\mathrm{Fl}\,R$ et $\mathrm{card}\,\mathrm{Fl}\,R'$

% slanted
\margin{dégager lemme sur syst. inductifs de catégories (qui servirait dans les 2 théorèmes...)}

% archive sheet number: 8
